\documentclass[12pt,a4paper]{article}
\usepackage[spanish]{babel}
\usepackage{geometry}
\usepackage{graphicx}
\usepackage{xcolor}
\usepackage{hyperref}
\usepackage{fancyhdr}
\usepackage{lastpage}
\usepackage{booktabs}
\usepackage{array}
\usepackage{colortbl}
\usepackage{multicol}
\usepackage{listings}
\usepackage{tikz}
\usepackage{enumerate}

\geometry{left=2.5cm, right=2.5cm, top=3cm, bottom=2.5cm}

\definecolor{primaryblue}{RGB}{0,102,204}
\definecolor{darkblue}{RGB}{0,51,102}
\definecolor{lightblue}{RGB}{230,242,255}
\definecolor{accentgreen}{RGB}{52,168,83}
\definecolor{darktext}{RGB}{51,51,51}

\hypersetup{
    colorlinks=true,
    linkcolor=primaryblue,
    urlcolor=primaryblue,
    citecolor=primaryblue
}

\pagestyle{fancy}
\fancyhf{}
\lhead{\textcolor{primaryblue}{\textbf{KAWSAYMI CARE}}}
\rhead{\textcolor{darktext}{Backend - NestJS}}
\cfoot{\textcolor{darktext}{Página \thepage\ de \pageref{LastPage}}}
\renewcommand{\headrulewidth}{1.5pt}
\renewcommand{\headrule}{\color{primaryblue}\hrule}

\usepackage{titlesec}
\titleformat{\section}{\Large\bfseries\color{primaryblue}}{}{0pt}{}[\color{primaryblue}\titlerule]
\titleformat{\subsection}{\large\bfseries\color{darkblue}}{}{0pt}{}
\titleformat{\subsubsection}{\normalsize\bfseries\color{darktext}}{}{0pt}{}

\title{
    \vspace{-2cm}
    \huge\textbf{\color{primaryblue}KAWSAYMI CARE}\\
    \Large\color{darkblue}Planificación de Construcción - Backend\\
    \normalsize\color{darktext}NestJS + Supabase + PostgreSQL
}
\author{Equipo de Desarrollo Backend}
\date{\today}

\begin{document}

\maketitle
\thispagestyle{fancy}

\vspace{1cm}

\begin{center}
    \boxed{\begin{minipage}{0.9\textwidth}
        \centering
        \textbf{\textcolor{primaryblue}{PLANIFICACIÓN DETALLADA DEL BACKEND}}\\[0.5cm]
        Este documento presenta la estructura, módulos, endpoints, base de datos\\
        y plan de construcción del backend con NestJS en 4 semanas.
    \end{minipage}}
\end{center}

\newpage
\tableofcontents
\newpage

% ============================================
\section{1. ARQUITECTURA DEL BACKEND}
% ============================================

\subsection{1.1 Stack Tecnológico}

\begin{itemize}
    \item \textbf{Framework:} NestJS 10.x con TypeScript
    \item \textbf{Base de Datos:} Supabase PostgreSQL
    \item \textbf{ORM:} Prisma
    \item \textbf{Autenticación:} Supabase Auth + JWT
    \item \textbf{Validación:} class-validator + class-transformer
    \item \textbf{Scheduler:} @nestjs/schedule (node-cron)
    \item \textbf{Logging:} @nestjs/common + Winston
    \item \textbf{Testing:} Jest
    \item \textbf{Deploy:} Railway / Fly.io
\end{itemize}

\subsection{1.2 Estructura de Carpetas}

\begin{verbatim}
src/
├── auth/
│   ├── auth.controller.ts
│   ├── auth.service.ts
│   ├── auth.module.ts
│   ├── strategies/
│   │   └── jwt.strategy.ts
│   └── guards/
│       └── jwt.guard.ts
├── users/
│   ├── users.controller.ts
│   ├── users.service.ts
│   ├── users.module.ts
│   └── dto/
│       ├── create-user.dto.ts
│       └── update-user.dto.ts
├── medications/
│   ├── medications.controller.ts
│   ├── medications.service.ts
│   ├── medications.module.ts
│   └── dto/
├── events/
│   ├── events.controller.ts
│   ├── events.service.ts
│   ├── events.module.ts
│   └── dto/
├── adherence/
│   ├── adherence.controller.ts
│   ├── adherence.service.ts
│   ├── adherence.module.ts
│   └── dto/
├── caregivers/
│   ├── caregivers.controller.ts
│   ├── caregivers.service.ts
│   ├── caregivers.module.ts
│   └── dto/
├── health/
│   ├── health.controller.ts
│   ├── health.service.ts
│   └── health.module.ts
├── scheduler/
│   ├── scheduler.service.ts
│   └── scheduler.module.ts
├── prisma/
│   └── prisma.service.ts
├── common/
│   ├── decorators/
│   ├── filters/
│   ├── interceptors/
│   └── pipes/
├── config/
│   └── database.config.ts
├── app.module.ts
└── main.ts
\end{verbatim}

\subsection{1.3 Principios de Arquitectura}

\begin{enumerate}
    \item \textbf{Modular:} Cada funcionalidad es un módulo independiente
    \item \textbf{Layered:} Separación clara entre controller, service, repository
    \item \textbf{Type-Safe:} TypeScript en todo el stack
    \item \textbf{Testeable:} Inyección de dependencias para facilitar testing
    \item \textbf{Scalable:} Preparado para crecer sin refactoring mayor
\end{enumerate}

% ============================================
\section{2. MÓDULOS DEL BACKEND}
% ============================================

\subsection{2.1 Auth Module}

\subsubsection{Responsabilidades}
\begin{itemize}
    \item Registro de usuarios (paciente/cuidador)
    \item Autenticación con Supabase Auth
    \item Generación de JWT tokens
    \item Refresh tokens
    \item Logout
\end{itemize}

\subsubsection{Endpoints}

\begin{center}
\begin{tabular}{|l|l|l|}
\hline
\rowcolor{lightblue}
\textbf{Método} & \textbf{Endpoint} & \textbf{Descripción} \\
\hline
POST & /auth/register & Registrar nuevo usuario \\
\hline
POST & /auth/login & Login con email y contraseña \\
\hline
POST & /auth/login-oauth & Login con OAuth (Google, Apple) \\
\hline
POST & /auth/refresh & Refrescar JWT token \\
\hline
POST & /auth/logout & Logout (invalidar token) \\
\hline
POST & /auth/forgot-password & Solicitar reset de contraseña \\
\hline
POST & /auth/reset-password & Confirmar reset de contraseña \\
\hline
\end{tabular}
\end{center}

\subsubsection{Modelos (Prisma)}

\begin{verbatim}
model User {
  id            String      @id @default(cuid())
  email         String      @unique
  role          Role
  name          String
  dateOfBirth   DateTime?
  location      String?
  language      String      @default("es")
  timezone      String      @default("UTC")
  allergies     String[]
  conditions    String[]
  createdAt     DateTime    @default(now())
  updatedAt     DateTime    @updatedAt
  
  medications   Medication[]
  events        MedicationEvent[]
  caregivers    CaregiverRelation[]
  health        HealthData?
}

enum Role {
  PATIENT
  CAREGIVER
}
\end{verbatim}

\subsection{2.2 Users Module}

\subsubsection{Responsabilidades}
\begin{itemize}
    \item CRUD de perfil de usuario
    \item Actualizar datos personales
    \item Gestionar alergias y condiciones
    \item Preferencias de usuario
\end{itemize}

\subsubsection{Endpoints}

\begin{center}
\begin{tabular}{|l|l|l|}
\hline
\rowcolor{lightblue}
\textbf{Método} & \textbf{Endpoint} & \textbf{Descripción} \\
\hline
GET & /users/me & Obtener perfil del usuario actual \\
\hline
PUT & /users/me & Actualizar perfil \\
\hline
PUT & /users/me/allergies & Actualizar alergias \\
\hline
PUT & /users/me/conditions & Actualizar condiciones médicas \\
\hline
DELETE & /users/me & Eliminar cuenta (soft delete) \\
\hline
\end{tabular}
\end{center}

\subsection{2.3 Medications Module}

\subsubsection{Responsabilidades}
\begin{itemize}
    \item CRUD de medicamentos
    \item Generación automática de eventos
    \item Gestión de estados (activo, completado, suspendido)
    \item Validación de dosis y frecuencia
\end{itemize}

\subsubsection{Endpoints}

\begin{center}
\begin{tabular}{|l|l|l|}
\hline
\rowcolor{lightblue}
\textbf{Método} & \textbf{Endpoint} & \textbf{Descripción} \\
\hline
GET & /medications & Listar medicamentos del usuario \\
\hline
GET & /medications/:id & Obtener detalles de medicamento \\
\hline
POST & /medications & Crear nuevo medicamento \\
\hline
PUT & /medications/:id & Actualizar medicamento \\
\hline
PATCH & /medications/:id/status & Cambiar estado (activo/suspendido) \\
\hline
DELETE & /medications/:id & Eliminar medicamento \\
\hline
\end{tabular}
\end{center}

\subsubsection{Modelos (Prisma)}

\begin{verbatim}
model Medication {
  id              String      @id @default(cuid())
  userId          String
  name            String
  dose            String
  frequency       Int
  intervalHours   Int
  instructions    String?
  startDate       DateTime
  endDate         DateTime?
  status          MedStatus   @default(ACTIVE)
  schedule        String[]
  createdAt       DateTime    @default(now())
  updatedAt       DateTime    @updatedAt
  
  user            User        @relation(fields: [userId], references: [id])
  events          MedicationEvent[]
}

enum MedStatus {
  ACTIVE
  COMPLETED
  SUSPENDED
}
\end{verbatim}

\subsection{2.4 Events Module}

\subsubsection{Responsabilidades}
\begin{itemize}
    \item Generación de eventos automáticos
    \item Actualización de estado (taken, missed, pending)
    \item Consulta de eventos por fecha/rango
\end{itemize}

\subsubsection{Endpoints}

\begin{center}
\begin{tabular}{|l|l|l|}
\hline
\rowcolor{lightblue}
\textbf{Método} & \textbf{Endpoint} & \textbf{Descripción} \\
\hline
GET & /events & Listar eventos (con filtros) \\
\hline
GET & /events/today & Eventos de hoy \\
\hline
GET & /events/week & Eventos de esta semana \\
\hline
PATCH & /events/:id/mark-taken & Marcar como tomado \\
\hline
PATCH & /events/:id/mark-missed & Marcar como omitido \\
\hline
\end{tabular}
\end{center}

\subsubsection{Modelo (Prisma)}

\begin{verbatim}
model MedicationEvent {
  id                String      @id @default(cuid())
  medicationId      String
  userId            String
  dateTimeScheduled DateTime
  status            EventStatus @default(PENDING)
  createdAt         DateTime    @default(now())
  updatedAt         DateTime    @updatedAt
  
  medication        Medication  @relation(fields: [medicationId], references: [id])
  user              User        @relation(fields: [userId], references: [id])
}

enum EventStatus {
  PENDING
  TAKEN
  MISSED
}
\end{verbatim}

\subsection{2.5 Adherence Module}

\subsubsection{Responsabilidades}
\begin{itemize}
    \item Cálculo de adherencia (%)
    \item Agregación por período (día, semana, mes)
    \item Alertas por baja adherencia
    \item Reportes longitudinales
\end{itemize}

\subsubsection{Endpoints}

\begin{center}
\begin{tabular}{|l|l|l|}
\hline
\rowcolor{lightblue}
\textbf{Método} & \textbf{Endpoint} & \textbf{Descripción} \\
\hline
GET & /adherence/today & Adherencia de hoy \\
\hline
GET & /adherence/week & Adherencia semanal \\
\hline
GET & /adherence/month & Adherencia mensual \\
\hline
GET & /adherence/stats & Estadísticas generales \\
\hline
\end{tabular}
\end{center}

\subsubsection{Cálculo de Adherencia}

\begin{verbatim}
adherence_percentage = (taken_events / total_events) * 100

Reglas:
- Si 2+ eventos consecutivos omitidos → ALERT
- Si adherence < 60% → ALERT
- Si adherence >= 80% → GOOD
\end{verbatim}

\subsection{2.6 Caregivers Module}

\subsubsection{Responsabilidades}
\begin{itemize}
    \item Crear relaciones paciente-cuidador
    \item Gestionar permisos (read, notify)
    \item Generar alertas para cuidadores
    \item Dashboard del cuidador
\end{itemize}

\subsubsection{Endpoints}

\begin{center}
\begin{tabular}{|l|l|l|}
\hline
\rowcolor{lightblue}
\textbf{Método} & \textbf{Endpoint} & \textbf{Descripción} \\
\hline
POST & /caregivers/invite & Invitar cuidador \\
\hline
GET & /caregivers/my-patients & Listar mis pacientes (cuidador) \\
\hline
GET & /caregivers/my-caregivers & Listar mis cuidadores (paciente) \\
\hline
PATCH & /caregivers/:id/permissions & Actualizar permisos \\
\hline
DELETE & /caregivers/:id & Eliminar relación \\
\hline
GET & /caregivers/:patientId/alerts & Alertas del paciente \\
\hline
\end{tabular}
\end{center}

\subsection{2.7 Health Module}

\subsubsection{Responsabilidades}
\begin{itemize}
    \item Gestionar peso, talla, IMC
    \item Hábitos: sueño, ejercicio, hidratación
    \item Detección de polifarmacia
    \item Interacciones medicamento-condición
\end{itemize}

\subsubsection{Endpoints}

\begin{center}
\begin{tabular}{|l|l|l|}
\hline
\rowcolor{lightblue}
\textbf{Método} & \textbf{Endpoint} & \textbf{Descripción} \\
\hline
GET & /health/profile & Obtener perfil de salud \\
\hline
POST & /health/weight & Registrar peso \\
\hline
GET & /health/imc & Calcular y obtener IMC \\
\hline
GET & /health/polypharmacy & Detectar polifarmacia \\
\hline
GET & /health/interactions & Verificar interacciones \\
\hline
\end{tabular}
\end{center}

\subsection{2.8 Scheduler Module}

\subsubsection{Responsabilidades}
\begin{itemize}
    \item Generar eventos automáticamente cada día
    \item Enviar notificaciones push en horarios
    \item Limpiar eventos antiguos
    \item Generar reportes periódicos
\end{itemize}

\subsubsection{Tareas Programadas}

\begin{verbatim}
// Cada día a las 00:00 UTC
- Generar eventos para próximos 30 días

// Cada medicamento a su hora
- Enviar notificación push (FCM)

// Cada semana (lunes 00:00 UTC)
- Generar reporte de adherencia

// Cada mes (1° a las 00:00 UTC)
- Generar estadísticas y trends
\end{verbatim}

% ============================================
\section{3. BASE DE DATOS - DIAGRAMA COMPLETO}
% ============================================

\subsection{3.1 Modelo Entidad-Relación}

\begin{verbatim}
┌─────────────────┐
│   User          │
├─────────────────┤
│ id (PK)         │
│ email           │
│ role            │
│ name            │
│ dateOfBirth     │
│ location        │
│ conditions[]    │
│ allergies[]     │
└─────────────────┘
        │
        ├──────────→ Medication (1:N)
        ├──────────→ MedicationEvent (1:N)
        ├──────────→ HealthData (1:1)
        └──────────→ CaregiverRelation (1:N)

┌──────────────────┐
│   Medication     │
├──────────────────┤
│ id (PK)          │
│ userId (FK)      │
│ name             │
│ dose             │
│ frequency        │
│ intervalHours    │
│ status           │
│ schedule[]       │
└──────────────────┘
        │
        └──────────→ MedicationEvent (1:N)

┌──────────────────────┐
│   MedicationEvent    │
├──────────────────────┤
│ id (PK)              │
│ medicationId (FK)    │
│ userId (FK)          │
│ dateTimeScheduled    │
│ status               │
└──────────────────────┘

┌────────────────────┐
│   HealthData       │
├────────────────────┤
│ id (PK)            │
│ userId (FK)        │
│ weight             │
│ height             │
│ imc                │
│ sleep_hours        │
│ exercise_minutes   │
│ water_liters       │
└────────────────────┘

┌────────────────────────┐
│   CaregiverRelation    │
├────────────────────────┤
│ patientId (FK)         │
│ caregiverId (FK)       │
│ permissions[]          │
│ createdAt              │
└────────────────────────┘
\end{verbatim}

% ============================================
\section{4. PLAN DE CONSTRUCCIÓN - 4 SEMANAS}
% ============================================

\subsection{4.1 SEMANA 1: Setup + Auth Module}

\subsubsection{Lunes-Martes: Infrastructure}
\begin{itemize}
    \item Crear proyecto NestJS
    \item Configurar Supabase y PostgreSQL
    \item Setup de Prisma
    \item Configurar variables de entorno
    \item CI/CD con GitHub Actions (básico)
\end{itemize}

\subsubsection{Miércoles-Viernes: Auth Module}
\begin{itemize}
    \item Crear AuthController y AuthService
    \item Implementar JWT Strategy
    \item Endpoints: register, login, logout
    \item Proteger rutas con @UseGuards(JwtAuthGuard)
    \item Testing unitario de Auth
\end{itemize}

\subsubsection{Entregables}
\begin{itemize}
    \item Proyecto con estructura base
    \item Auth completamente funcional
    \item Usuarios pueden registrarse y loguearse
\end{itemize}

\subsection{4.2 SEMANA 2: Users + Medications Module}

\subsubsection{Lunes-Martes: Users Module}
\begin{itemize}
    \item CRUD de perfil de usuario
    \item Gestión de alergias y condiciones
    \item Endpoints GET/PUT /users/me
    \item Validaciones de datos
    \item Tests unitarios
\end{itemize}

\subsubsection{Miércoles-Viernes: Medications Module}
\begin{itemize}
    \item CRUD de medicamentos
    \item Validación de dosis y frecuencia
    \item Generación automática de eventos
    \item Endpoints de medicamentos
    \item Tests unitarios
\end{itemize}

\subsubsection{Entregables}
\begin{itemize}
    \item Usuarios pueden actualizar su perfil
    \item CRUD de medicamentos funcional
    \item Eventos se generan automáticamente
\end{itemize}

\subsection{4.3 SEMANA 3: Events + Adherence + Health}

\subsubsection{Lunes: Events Module}
\begin{itemize}
    \item CRUD de eventos
    \item Mark as taken/missed
    \item Filtros por fecha
    \item Tests
\end{itemize}

\subsubsection{Martes-Miércoles: Adherence Module}
\begin{itemize}
    \item Cálculo de adherencia
    \item Agregación por período
    \item Alertas automáticas
    \item Endpoint de estadísticas
    \item Tests
\end{itemize}

\subsubsection{Jueves-Viernes: Health Module}
\begin{itemize}
    \item CRUD de health data
    \item Cálculo de IMC
    \item Detección de polifarmacia
    \item Verificación de interacciones
    \item Tests
\end{itemize}

\subsubsection{Entregables}
\begin{itemize}
    \item Sistema de eventos completo
    \item Cálculo de adherencia funcionando
    \item Health data y análisis
\end{itemize}

\subsection{4.4 SEMANA 4: Caregivers + Scheduler + Testing}

\subsubsection{Lunes-Martes: Caregivers Module}
\begin{itemize}
    \item Crear relaciones paciente-cuidador
    \item Gestionar permisos
    \item Dashboard del cuidador
    \item Alertas para cuidadores
    \item Tests
\end{itemize}

\subsubsection{Miércoles: Scheduler Module}
\begin{itemize}
    \item Setup de node-cron
    \item Generación de eventos automática
    \item Scheduler de notificaciones
    \item Limpieza de datos antiguos
    \item Tests
\end{itemize}

\subsubsection{Jueves-Viernes: Testing + Polish}
\begin{itemize}
    \item Testing E2E de flujos principales
    \item Integration tests
    \item Bug fixes
    \item Documentación con Swagger
    \item Deploy a staging
\end{itemize}

\subsubsection{Entregables}
\begin{itemize}
    \item Backend completamente funcional
    \item Todos los módulos integrados
    \item Testing > 80\% coverage
    \item Documentación API (Swagger)
    \item Listo para producción
\end{itemize}

% ============================================
\section{5. ENDPOINTS COMPLETOS}
% ============================================

\subsection{5.1 Auth Endpoints}

\begin{verbatim}
POST /auth/register
Body: { email, password, name, role }
Response: { user, token, refreshToken }

POST /auth/login
Body: { email, password }
Response: { user, token, refreshToken }

POST /auth/refresh
Headers: { Authorization: Bearer <token> }
Response: { token, refreshToken }

POST /auth/logout
Headers: { Authorization: Bearer <token> }
Response: { success: true }
\end{verbatim}

\subsection{5.2 Users Endpoints}

\begin{verbatim}
GET /users/me
Headers: { Authorization: Bearer <token> }
Response: { id, email, name, role, ... }

PUT /users/me
Body: { name, location, timezone, language }
Response: { updated user }

PUT /users/me/conditions
Body: { conditions: [string] }
Response: { conditions }

PUT /users/me/allergies
Body: { allergies: [string] }
Response: { allergies }
\end{verbatim}

\subsection{5.3 Medications Endpoints}

\begin{verbatim}
GET /medications
Query: { status?, skip?, take? }
Response: [{ id, name, dose, frequency, ... }]

POST /medications
Body: { name, dose, frequency, intervalHours, schedule, ... }
Response: { created medication }

PUT /medications/:id
Body: { name?, dose?, frequency?, ... }
Response: { updated medication }

PATCH /medications/:id/status
Body: { status: "ACTIVE" | "SUSPENDED" | "COMPLETED" }
Response: { medication with new status }

DELETE /medications/:id
Response: { success: true }
\end{verbatim}

\subsection{5.4 Events Endpoints}

\begin{verbatim}
GET /events
Query: { dateFrom?, dateTo?, medicationId?, status? }
Response: [{ id, status, dateTimeScheduled, ... }]

GET /events/today
Response: [events for today]

GET /events/week
Response: [events for this week]

PATCH /events/:id/mark-taken
Body: { timeMarked?: DateTime }
Response: { event with status=TAKEN }

PATCH /events/:id/mark-missed
Response: { event with status=MISSED }
\end{verbatim}

% ============================================
\section{6. CONVENCIONES Y ESTÁNDARES}
% ============================================

\subsection{6.1 Nombres}

\begin{itemize}
    \item \textbf{Controllers:} PascalCase + Controller
    \item \textbf{Services:} PascalCase + Service
    \item \textbf{Modules:} PascalCase + Module
    \item \textbf{DTOs:} PascalCase + Dto
    \item \textbf{Métodos:} camelCase
    \item \textbf{Variables:} camelCase
\end{itemize}

\subsection{6.2 Validaciones}

\begin{verbatim}
// DTOs deben tener validaciones
import { IsEmail, IsNotEmpty, IsString } from 'class-validator';

export class CreateUserDto {
  @IsEmail()
  email: string;

  @IsNotEmpty()
  @IsString()
  password: string;

  @IsNotEmpty()
  @IsString()
  name: string;
}
\end{verbatim}

\subsection{6.3 Manejo de Errores}

\begin{verbatim}
// Usar excepciones de NestJS
throw new BadRequestException('Invalid input');
throw new UnauthorizedException('Not authenticated');
throw new NotFoundException('Resource not found');
throw new ConflictException('Resource already exists');
throw new InternalServerErrorException('Server error');
\end{verbatim}

\subsection{6.4 Logging}

\begin{verbatim}
// Usar logger de NestJS
private readonly logger = new Logger(MedicationsService.name);

this.logger.log('Medication created: ' + id);
this.logger.error('Error creating medication', error);
this.logger.warn('Low adherence detected');
\end{verbatim}

% ============================================
\section{7. TESTING STRATEGY}
% ============================================

\subsection{7.1 Unit Tests}

\begin{itemize}
    \item Cada service debe tener tests
    \item Mock de dependencias con jest.mock()
    \item Coverage mínimo: 80\%
    \item Archivo: \texttt{service.spec.ts}
\end{itemize}

\subsection{7.2 Integration Tests}

\begin{itemize}
    \item Pruebas E2E de flujos principales
    \item Usar TestingModule de NestJS
    \item Pruebas con BD real (test database)
    \item Coverage de endpoints críticos
\end{itemize}

\subsection{7.3 Ejecución de Tests}

\begin{verbatim}
# Tests unitarios
npm run test

# Tests E2E
npm run test:e2e

# Coverage
npm run test:cov
\end{verbatim}

% ============================================
\section{8. DEPLOYMENT Y PRODUCCIÓN}
% ============================================

\subsection{8.1 Preparación}

\begin{itemize}
    \item Build del proyecto: \texttt{npm run build}
    \item Variables de entorno en producción
    \item Logs y monitoreo configurados
    \item Database backups automáticos en Supabase
\end{itemize}

\subsection{8.2 Deploy a Railway/Fly.io}

\begin{verbatim}
# Railway
railway link
railway up

# Fly.io
flyctl launch
flyctl deploy
\end{verbatim}

\subsection{8.3 Monitoreo}

\begin{itemize}
    \item \textbf{Logs:} Supabase logs + Winston
    \item \textbf{Errores:} Sentry integration
    \item \textbf{Performance:} Railway/Fly.io dashboards
    \item \textbf{Uptime:} Status page
\end{itemize}

% ============================================
\section{9. SEGURIDAD}
% ============================================

\subsection{9.1 Autenticación}

\begin{itemize}
    \item JWT tokens firmados con secret
    \item Refresh tokens con expiración
    \item Supabase Auth para OAuth2
\end{itemize}

\subsection{9.2 Autorización}

\begin{itemize}
    \item Guards para proteger endpoints
    \item RLS (Row Level Security) en Supabase
    \item Validar que usuario solo acceda sus datos
\end{itemize}

\subsection{9.3 Datos Sensibles}

\begin{itemize}
    \item Hashing de contraseñas con bcrypt (manejado por Supabase)
    \item Encriptación en tránsito (HTTPS)
    \item Encriptación en reposo en Supabase
\end{itemize}

% ============================================
\section{10. DOCUMENTACIÓN}
% ============================================

\subsection{10.1 Swagger/OpenAPI}

\begin{verbatim}
import { NestFactory } = from '@nestjs/core';
import { SwaggerModule, DocumentBuilder } from '@nestjs/swagger';
import { AppModule } from './app.module';

async function bootstrap() {
  const app = await NestFactory.create(AppModule);
  
  const config = new DocumentBuilder()
    .setTitle('KAWSAYMI CARE API')
    .setDescription('API Backend')
    .setVersion('1.0')
    .addBearerAuth()
    .build();
    
  const document = SwaggerModule.createDocument(app, config);
  SwaggerModule.setup('api/docs', app, document);

  await app.listen(3000);
}

bootstrap();
\end{verbatim}

\subsection{10.2 README}

\begin{itemize}
    \item Guía de instalación
    \item Variables de entorno (.env.example)
    \item Cómo ejecutar en local
    \item Estructura de carpetas
    \item Cómo contribuir
\end{itemize}

\vspace{2cm}

\begin{center}
    \boxed{\begin{minipage}{0.8\textwidth}
        \centering
        \textbf{\textcolor{primaryblue}{Backend sólido, escalable y mantenible.}}\\
        \vspace{0.3cm}
        \textit{Con NestJS + Supabase, el MVP está preparado\\
        para escalar a millones de eventos.}
    \end{minipage}}
\end{center}

\end{document}